% ============================================================
\chapter{Chaos --- D\'efinitions et Propri\'et\'es}
\label{ch:chaos}
% ============================================================

\section{Introduction}

En 1963, Edward Lorenz, m\'et\'eorologue au MIT, d\'ecouvre par accident que son mod\`ele simplifi\'e de convection atmosph\'erique produit des r\'esultats radicalement diff\'erents selon que l'on arrondit une condition initiale \`a trois ou six d\'ecimales. Ce ph\'enom\`ene --- la \emph{sensibilit\'e aux conditions initiales} --- deviendra le symb\^ole du chaos d\'eterministe et sera popularis\'e sous le nom d'\og effet papillon \fg. Mais qu'est-ce que le chaos, exactement~? La question est plus subtile qu'il n'y para\^it. Robert Devaney proposera en 1989 une d\'efinition rigoureuse combinant trois ingr\'edients~: sensibilit\'e aux conditions initiales, transitivit\'e topologique et densit\'e des orbites p\'eriodiques. Li et Yorke, dans leur article c\'el\`ebre \og Period Three Implies Chaos \fg{} (1975), en avaient donn\'e une autre. Ce chapitre explore ces d\'efinitions, leurs relations, et les propri\'et\'es fondamentales des syst\`emes chaotiques --- ces syst\`emes d\'eterministes qui, paradoxalement, produisent l'apparence du hasard.

\section{Sensibilit\'e aux conditions initiales}

\begin{definition}[Sensibilit\'e aux conditions initiales]
Un syst\`eme dynamique $\dot{x} = f(x)$ (ou $x_{n+1} = f(x_n)$) poss\`ede
une \textbf{sensibilit\'e aux conditions initiales} (SCI) sur un ensemble
invariant $\Lambda$ s'il existe $\delta > 0$ tel que pour tout $x \in \Lambda$
et tout $\epsilon > 0$, il existe $y$ avec $\norm{y - x} < \epsilon$ et
$t > 0$ (ou $n > 0$) tel que
\[
    \norm{\varphi(t, x) - \varphi(t, y)} > \delta.
\]
\end{definition}

\begin{intuition}
La sensibilit\'e aux conditions initiales est souvent appel\'ee
\og effet papillon \fg : une perturbation infime peut conduire \`a des
\'evolutions radicalement diff\'erentes. C'est la d\'ecouverte fondamentale
de Lorenz (1963).
\end{intuition}

\begin{remarque}
La sensibilit\'e seule ne suffit pas \`a d\'efinir le chaos. L'application
$x \mapsto 2x$ sur $\R$ est sensible aux conditions initiales mais n'est
pas chaotique (les orbites s'\'echappent simplement vers l'infini).
\end{remarque}

\section{D\'efinition de Devaney du chaos}

\begin{definition}[Chaos au sens de Devaney]
\label{def:devaney}
Un syst\`eme dynamique $f : X \to X$ sur un espace m\'etrique $(X, d)$ est
\textbf{chaotique au sens de Devaney} s'il satisfait les trois conditions :
\begin{enumerate}
    \item \textbf{Transitivit\'e topologique} : pour tous ouverts non vides
    $U, V \subseteq X$, il existe $n \in \N$ tel que $f^n(U) \cap V \neq \emptyset$.
    \item \textbf{Densit\'e des orbites p\'eriodiques} : l'ensemble des points
    p\'eriodiques est dense dans $X$.
    \item \textbf{Sensibilit\'e aux conditions initiales}.
\end{enumerate}
\end{definition}

\begin{theoreme}[Banks et al., 1992]
Si $f : X \to X$ est transitive topologiquement et a des orbites p\'eriodiques
denses sur un espace m\'etrique infini, alors $f$ poss\`ede automatiquement
la sensibilit\'e aux conditions initiales. Ainsi, la condition (3) est
\textbf{redondante} dans la d\'efinition de Devaney.
\end{theoreme}

\section{D\'efinition de Li-Yorke du chaos}

\begin{definition}[Paire de Li-Yorke]
Deux points $x, y$ forment une \textbf{paire de Li-Yorke} si
\[
    \limsup_{n \to \infty} \norm{f^n(x) - f^n(y)} > 0 \qquad \text{et} \qquad
    \liminf_{n \to \infty} \norm{f^n(x) - f^n(y)} = 0.
\]
\end{definition}

\begin{definition}[Chaos au sens de Li-Yorke]
Le syst\`eme est \textbf{chaotique au sens de Li-Yorke} s'il existe un
ensemble non d\'enombrable $S$ (ensemble brouill\'e) tel que toute paire
de points distincts de $S$ est une paire de Li-Yorke.
\end{definition}

\begin{theoreme}[Li-Yorke, 1975]
\label{thm:li-yorke}
Soit $f : [0,1] \to [0,1]$ continue. Si $f$ poss\`ede une orbite de
p\'eriode 3, alors $f$ poss\`ede des orbites de toute p\'eriode $n \in \N^*$,
et $f$ est chaotique au sens de Li-Yorke.
\end{theoreme}

\begin{remarque}
L'\'enonc\'e populaire \og Period three implies chaos \fg r\'esume ce
th\'eor\`eme. Plus g\'en\'eralement, l'ordre de Sharkovskii d\'etermine
quelles p\'eriodes impliquent quelles autres.
\end{remarque}

\section{L'ordre de Sharkovskii}

\begin{theoreme}[Sharkovskii, 1964]
D\'efinissons l'ordre de Sharkovskii sur $\N^*$ :
\[
    3 \triangleright 5 \triangleright 7 \triangleright \cdots \triangleright
    2\cdot 3 \triangleright 2\cdot 5 \triangleright \cdots \triangleright
    4\cdot 3 \triangleright 4\cdot 5 \triangleright \cdots \triangleright
    \cdots \triangleright 8 \triangleright 4 \triangleright 2 \triangleright 1.
\]
Si $f : \R \to \R$ est continue et a une orbite de p\'eriode $m$, alors $f$
a des orbites de toute p\'eriode $n$ telle que $m \triangleright n$.
\end{theoreme}

\section{Exemples fondamentaux}

\subsection{L'application de d\'ecalage (shift)}

\begin{definition}[Espace des suites et shift]
L'espace des suites binaires est
$\Sigma_2 = \{0, 1\}^\N = \{s = (s_0, s_1, s_2, \ldots) : s_i \in \{0, 1\}\}$
muni de la m\'etrique
$d(s, t) = \sum_{i=0}^{\infty} \frac{\abs{s_i - t_i}}{2^i}$.
Le \textbf{shift} $\sigma : \Sigma_2 \to \Sigma_2$ est d\'efini par
$\sigma(s_0, s_1, s_2, \ldots) = (s_1, s_2, s_3, \ldots)$.
\end{definition}

\begin{theoreme}
Le shift $\sigma$ sur $\Sigma_2$ est chaotique au sens de Devaney :
\begin{enumerate}
    \item Il est topologiquement transitif.
    \item Les points p\'eriodiques sont denses.
    \item Il est sensible aux conditions initiales avec $\delta = 1$.
\end{enumerate}
\end{theoreme}

\subsection{L'application de la tente}

\begin{definition}[Application de la tente]
L'\textbf{application de la tente} est $T : [0,1] \to [0,1]$ d\'efinie par
\[
    T(x) = \begin{cases} 2x & \text{si } 0 \leq x \leq 1/2, \\ 2(1-x) & \text{si } 1/2 < x \leq 1. \end{cases}
\]
\end{definition}

\begin{proposition}
L'application de la tente est semi-conjugu\'ee au shift sur $\Sigma_2$
et est chaotique au sens de Devaney.
\end{proposition}

\subsection{L'application logistique \`a $r = 4$}

\begin{theoreme}
L'application logistique $f_4(x) = 4x(1-x)$ sur $[0,1]$ est topologiquement
conjugu\'ee \`a l'application de la tente via le changement de variable
$h(x) = \frac{2}{\pi}\arcsin(\sqrt{x})$. Elle est donc chaotique au sens
de Devaney.
\end{theoreme}

\section{Le fer \`a cheval de Smale}

\begin{definition}[Fer \`a cheval de Smale]
Le \textbf{fer \`a cheval de Smale} est un diff\'eomorphisme du plan qui
\'etire un carr\'e, le plie en forme de fer \`a cheval, et le remet dans
le carr\'e initial. L'ensemble invariant maximal est un ensemble de Cantor
sur lequel la dynamique est conjugu\'ee au shift sur $\Sigma_2$.
\end{definition}

\begin{theoreme}[Smale, 1967]
Sur l'ensemble invariant du fer \`a cheval, la dynamique est :
\begin{enumerate}
    \item chaotique au sens de Devaney ;
    \item hyperbolique (uniformly hyperbolic) ;
    \item structurellement stable (persiste sous petites perturbations).
\end{enumerate}
\end{theoreme}

\begin{remarque}
Le fer \`a cheval de Smale est le prototype du chaos \textbf{hyperbolique}.
Il appara\^it g\'en\'eriquement pr\`es des orbites homoclines transverses
(th\'eor\`eme de Smale-Birkhoff).
\end{remarque}

\section{Routes vers le chaos}

\begin{definition}[Cascade de doublements de p\'eriode]
La \textbf{cascade de Feigenbaum} est une suite infinie de bifurcations de
doublement de p\'eriode $1 \to 2 \to 4 \to 8 \to \cdots$ menant au chaos.
\end{definition}

\begin{theoreme}[Universalit\'e de Feigenbaum]
Les valeurs de bifurcation $\mu_n$ satisfont
\[
    \lim_{n \to \infty} \frac{\mu_n - \mu_{n-1}}{\mu_{n+1} - \mu_n} = \delta_F \approx 4{,}6692\ldots
\]
La constante $\delta_F$ est \textbf{universelle} : elle ne d\'epend pas de
la famille d'applications consid\'er\'ee.
\end{theoreme}

\begin{formules}[title=\textbf{Constantes de Feigenbaum}]
\begin{align*}
    \delta_F &\approx 4{,}669\,201\,609\,\ldots \quad \text{(rapport des intervalles de param\`etre)} \\
    \alpha_F &\approx -2{,}502\,907\,875\,\ldots \quad \text{(facteur d'\'echelle spatial)}
\end{align*}
\end{formules}

Autres routes vers le chaos :
\begin{itemize}
    \item \textbf{Intermittence} (Pomeau-Manneville) : alternance de phases
    laminaires et de bursts chaotiques.
    \item \textbf{Quasi-p\'eriodicit\'e} (Ruelle-Takens) : passage par un
    tore invariant qui se d\'etruit.
    \item \textbf{Crise} : collision d'un attracteur chaotique avec un
    r\'epulseur.
\end{itemize}

\section{Chaos dans les syst\`emes continus}

\begin{attention}[title=\textbf{Dimension minimale pour le chaos continu}]
Un syst\`eme autonome continu $\dot{x} = f(x)$ ne peut \^etre chaotique
qu'en dimension $n \geq 3$. En dimension 2, le th\'eor\`eme de
Poincar\'e-Bendixson interdit le chaos.
\end{attention}

\begin{exemple}[Th\'eor\`eme de Mel'nikov]
Le th\'eor\`eme de Mel'nikov permet de d\'etecter le chaos dans les syst\`emes
hamiltoniens perturb\'es. Pour le pendule forc\'e
$\ddot{x} + \sin x = \epsilon(\gamma\cos\omega t - \delta\dot{x})$,
la fonction de Mel'nikov pr\'edit l'existence d'orbites homoclines
transverses (et donc d'un fer \`a cheval de Smale) pour $\epsilon$
suffisamment petit.
\end{exemple}

\begin{codeblock}[title=\textbf{D\'emonstration de la sensibilit\'e aux conditions initiales}]
\begin{minted}{python}
import numpy as np
from scipy.integrate import solve_ivp
import matplotlib.pyplot as plt

def lorenz(t, state, sigma=10, rho=28, beta=8/3):
    x, y, z = state
    return [sigma*(y - x), rho*x - y - x*z, x*y - beta*z]

# Two nearby initial conditions
x0_1 = [1.0, 1.0, 1.0]
x0_2 = [1.0 + 1e-10, 1.0, 1.0]

t_span = [0, 50]
t_eval = np.linspace(0, 50, 5000)
sol1 = solve_ivp(lorenz, t_span, x0_1, t_eval=t_eval, max_step=0.01)
sol2 = solve_ivp(lorenz, t_span, x0_2, t_eval=t_eval, max_step=0.01)

fig, (ax1, ax2) = plt.subplots(2, 1, figsize=(12, 8))
ax1.plot(sol1.t, sol1.y[0], 'b-', lw=0.5, label='CI 1')
ax1.plot(sol2.t, sol2.y[0], 'r-', lw=0.5, label='CI 2')
ax1.set_xlabel('$t$'); ax1.set_ylabel('$x(t)$')
ax1.legend(); ax1.set_title('Sensibilite aux conditions initiales')

diff = np.sqrt(np.sum((sol1.y - sol2.y)**2, axis=0))
ax2.semilogy(sol1.t, diff, 'k-', lw=0.5)
ax2.set_xlabel('$t$'); ax2.set_ylabel('$\\|\\Delta x(t)\\|$')
ax2.set_title('Divergence exponentielle')
plt.tight_layout()
plt.savefig("sensitivity.pdf")
plt.show()
\end{minted}
\end{codeblock}

\section{Exercices}

\begin{exercice}[Shift et chaos]
Montrer que le shift $\sigma$ sur $\Sigma_2$ est topologiquement transitif
en construisant explicitement une orbite dense.
\end{exercice}

\begin{exercice}[P\'eriode 3]
Montrer que l'application $f(x) = x^2 - 7/4$ sur $\R$ a une orbite de
p\'eriode 3 et en d\'eduire qu'elle est chaotique au sens de Li-Yorke.
\end{exercice}

\begin{exercice}[Application de la tente]
Trouver tous les points p\'eriodiques de p\'eriode $\leq 3$ de l'application
de la tente et v\'erifier qu'ils sont instables.
\end{exercice}

\begin{exercice}[Logistique]
Pour l'application logistique $f_r(x) = rx(1-x)$, montrer num\'eriquement
la cascade de doublement de p\'eriode et estimer la constante de Feigenbaum.
\end{exercice}

\begin{exercice}[Fer \`a cheval]
D\'ecrire g\'eom\'etriquement la construction du fer \`a cheval de Smale et
montrer que l'ensemble invariant est un produit de Cantor.
\end{exercice}

\begin{exercice}[Chaos en dimension 3]
Pourquoi le syst\`eme $\dot{x} = -y - z$, $\dot{y} = x + ay$,
$\dot{z} = b + z(x-c)$ (R\"ossler) peut-il \^etre chaotique alors qu'il
n'a que trois \'equations ? V\'erifier num\'eriquement pour $a = 0.2$,
$b = 0.2$, $c = 5.7$.
\end{exercice}
